\section{Changing one adjacent order}\label{changing-one-adjacent-order}

Let \((B;a,b,c)\) be a rank-three string C-group of type \(\{h,q\}\), with
\(h,q\) even, and suppose the map

\[\rho:B\longrightarrow B_{12},\qquad a\mapsto1, b\mapsto b, c\mapsto c
\tag{7}\]

is a retraction. For \(k\) a positive multiple of \(q\), take the standard
marked epimorphism \(D_{2k}\to B_{12}=D_{2q}\). The compatible product

\[F=\{(g,d)\in B\times D_{2k}:\rho(g)=\overline d\}\tag{8}\]

has markers \((a,1),(b,\widetilde b),(c,\widetilde c)\). The evenness of
\(h\) is forced by (7); it also ensures that the first pair has unchanged
order \(h\).

\textbf{Theorem AL.4 (one-edge extension).} These markers generate \(F\). If \(B\) is finite, its order
is \(|B|k/q\). Its type is \(\{h,k\}\), and it has the intersection
property. The projection onto \(B\) is injective on the first adjacent
subgroup. If \(B\) is solvable, so is \(F\).

\emph{Proof.} The word \((r_1r_2)^q\) generates the entire cyclic kernel of
(8), since its first coordinate is trivial and its second has order
\(k/q\). Surjectivity onto \(B\) then proves generation and the order formula.
The first adjacent subgroup has second coordinate of order at most two,
with that coordinate determined by the image under (7); thus its
projection is injective. The second adjacent subgroup is the full
\(D_{2k}\). If an element belongs to both, its projection belongs to
\(\langle b\rangle\) by IP in \(B\). Subtract the appropriate common
involution. The remaining element lies in the first subgroup and in the
projection kernel, hence is trivial. This proves IP by the rank-three
criterion. Solvability follows from the cyclic kernel. \(\square\)

Theorem AL.4 is a specialization of UA.SK; it is included to make clear that
only the indicated deletion is needed here. More simultaneous edges use
the supported coordinate kernels of that theorem.

There is a useful exact presentation variant. Suppose a complete marked
presentation for \(B\) contains \((bc)^q=1\). Replace that relation by
\((bc)^k=1\), retain the remaining relators, and add
\([a,(bc)^q]=1\). Call the resulting presented group \(H\). If every retained
relator holds in (8), then \(H=F\) as a marked group. Indeed,
\(C=\langle(bc)^q\rangle\) is normal in \(H\): the last relation centralizes
it by \(a\), and \(b,c\) invert it. Its order is at most \(k/q\), and \(H/C=B\)
by the presentation. Consequently \(|H|\leq|B|k/q\). The marked
surjection onto (8) attains this bound, proving equality and injectivity.
The same argument is valid for an infinite base if the finite kernel
maps injectively and the induced map on the base is an isomorphism.
